\section{Public FreshQA Robustness Track: Retrieval Dominance And Calibration}

The public FreshQA-derived benchmark is the most important journal extension because it tests the controller on a public, time-varying source rather than only on the bundled internal track. The result is mixed but informative.

\begin{table}[h]
\centering
\small
\begin{tabular}{llrrrrr}
\toprule
Model & Mode & Quality & 95\% CI & Retrieve & Adapt & Stale \\
\midrule
flan-t5-base & always\_retrieve & 1.000 & [1.000, 1.000] & 192.0 & 0.0 & 0.000 \\
flan-t5-base & router & 0.582 & [0.442, 0.721] & 102.0 & 32.3 & 0.009 \\
flan-t5-base & router\_calibrated & 0.842 & [0.811, 0.873] & 155.7 & 4.0 & 0.000 \\
flan-t5-base & router\_no\_v\_future & 0.969 & [0.969, 0.969] & 181.0 & 4.0 & 0.000 \\
flan-t5-small & always\_retrieve & 1.000 & [1.000, 1.000] & 192.0 & 0.0 & 0.000 \\
flan-t5-small & router & 0.564 & [0.409, 0.719] & 101.7 & 32.0 & 0.003 \\
flan-t5-small & router\_calibrated & 0.844 & [0.819, 0.869] & 156.0 & 4.0 & 0.000 \\
flan-t5-small & router\_no\_v\_future & 0.964 & [0.964, 0.964] & 181.0 & 4.0 & 0.000 \\
\bottomrule
\end{tabular}

\caption{Public FreshQA-derived main slice from \texttt{artifacts/journal\_freshqa\_main\_calibrated\_v1/aggregate\_table.md}.}
\label{tab:j-fresh-main}
\end{table}

On the public main slice, retrieval remains dominant, but the story is now more complete. \texttt{always\_retrieve} achieves perfect quality on both backbones. The full router is substantially worse (`0.582` on \texttt{flan-t5-base}, `0.564` on \texttt{flan-t5-small}), which preserves the original diagnosis that delayed utility is too eager for a retrieval-dominant public regime. The new calibrated router materially improves that result without changing the action set: it rises to `0.842` on \texttt{flan-t5-base} and `0.844` on \texttt{flan-t5-small}. The no-$V_{\text{future}}$ ablation still performs best (`0.969` and `0.964`), so the correct interpretation is not that calibration solves the public benchmark completely. It is that regime-aware future-value gating is the right next mechanism, and that the public-track weakness is genuinely a delayed-utility problem rather than a generic failure of memory or adaptation.

\begin{table}[h]
\centering
\small
\begin{tabular}{llrrrrr}
\toprule
Model & Mode & Quality & 95\% CI & Retrieve & Adapt & Forgetting $\Delta$ \\
\midrule
flan-t5-base & always\_retrieve & 1.000 & [1.000, 1.000] & 384.0 & 0.0 & 0.000 \\
flan-t5-base & router & 0.751 & [0.676, 0.825] & 282.0 & 64.7 & -0.345 \\
flan-t5-base & router\_calibrated & 0.766 & [0.749, 0.783] & 289.3 & 4.0 & -0.385 \\
flan-t5-base & router\_no\_v\_future & 0.977 & [0.977, 0.977] & 374.0 & 4.0 & 0.000 \\
flan-t5-small & always\_retrieve & 1.000 & [1.000, 1.000] & 384.0 & 0.0 & 0.000 \\
flan-t5-small & router & 0.740 & [0.678, 0.803] & 281.3 & 64.0 & -0.345 \\
flan-t5-small & router\_calibrated & 0.787 & [0.753, 0.822] & 299.0 & 4.0 & -0.317 \\
flan-t5-small & router\_no\_v\_future & 0.974 & [0.974, 0.974] & 373.0 & 4.0 & -0.012 \\
\bottomrule
\end{tabular}

\caption{Public FreshQA-derived sequence slice from \texttt{artifacts/journal\_freshqa\_sequence\_calibrated\_v1/aggregate\_table.md}.}
\label{tab:j-fresh-seq}
\end{table}

The sequence slice sharpens the diagnosis. Calibration helps there too, but more modestly. The full router reaches `0.751` on \texttt{flan-t5-base} and `0.740` on \texttt{flan-t5-small}; the calibrated router improves those to `0.766` and `0.787`. However, the no-$V_{\text{future}}$ ablation again outperforms both on quality and forgetting. The sequence slice therefore becomes a main experiment rather than a supporting table: it shows that \texttt{FAST\_ADAPT} remains useful, but that future value must still be applied much more conservatively on public changing-answer sequences than on recurring-memory or bundled stable-confirmation regimes.

\begin{table}[h]
\centering
\small
\begin{tabular}{lrrrrr}
\toprule
Model / scale & Quality & Retrieve & Adapt & Stale & Forgetting $\Delta$ \\
\midrule
base / 0.00 & 0.932 & 171.0 & 4.0 & 0.000 & -0.094 \\
base / 0.50 & 0.670 & 119.3 & 35.7 & 0.007 & -0.427 \\
base / 1.00 & 0.580 & 101.7 & 32.0 & 0.009 & -0.615 \\
small / 0.00 & 0.958 & 180.0 & 4.0 & 0.000 & -0.031 \\
small / 0.50 & 0.656 & 119.7 & 36.3 & 0.002 & -0.427 \\
small / 1.00 & 0.562 & 101.3 & 31.7 & 0.003 & -0.615 \\
\bottomrule
\end{tabular}

\caption{Public FreshQA delayed-utility calibration sweep from \texttt{artifacts/journal\_freshqa\_future\_scale\_sweep\_v1/aggregate\_table.md}.}
\label{tab:j-fresh-scale}
\end{table}

Table~\ref{tab:j-fresh-scale} turns that diagnosis into a concrete next step. The best observed public-track quality occurs at \texttt{future\_value\_scale = 0.0} on both models. As the delayed-utility scale increases, quality and forgetting degrade monotonically. This is why the next controller step should be calibration, not a broader redesign: the public track is telling us that future value must be used far more selectively for changing-answer benchmarks than for recurring-memory benchmarks.

The cross-family robustness check tells the same story. On \texttt{SmolLM2-360M-Instruct}, \texttt{always\_retrieve} again reaches `1.000`, the full router drops to `0.569`, and the calibrated router recovers to `0.830`. The public boundary result and the value of regime-aware calibration are therefore not limited to the Flan family. We back that interpretation with a public provenance bundle and a paired manual audit bundle for the public track.
